% RDTII Rocky - Setup guide for a new computer.
% Plain LaTeX, pdfLaTeX, Overleaf-ready: upload this one file and press Recompile.
\documentclass[11pt,a4paper]{article}

\usepackage[utf8]{inputenc}
\usepackage[T1]{fontenc}
\usepackage{lmodern}   % vector fonts under T1 on every installation (a basic MiKTeX has only bitmap EC fonts)
\usepackage[margin=2.3cm]{geometry}
\usepackage{xcolor}
\usepackage{listings}
\usepackage{booktabs}
\usepackage{enumitem}
\usepackage{parskip}
\usepackage[colorlinks=true,linkcolor=rdtiiblue,urlcolor=rdtiiblue]{hyperref}

\definecolor{rdtiiblue}{HTML}{1A4D8F}
\definecolor{shell}{HTML}{F4F6F8}
\definecolor{rule}{HTML}{C9D3DD}

\lstset{
  basicstyle=\ttfamily\small,
  backgroundcolor=\color{shell},
  frame=single,
  rulecolor=\color{rule},
  framesep=6pt,
  breaklines=true,
  columns=fullflexible,
  keepspaces=true,
  showstringspaces=false,
  aboveskip=8pt, belowskip=8pt,
}

\newcommand{\win}{\textbf{Windows}}
\newcommand{\mac}{\textbf{Mac}}

\title{\textbf{RDTII Rocky}\\[2mm]
\large Setup guide: getting the tool running on a computer that has nothing installed}
\author{Team: Rocky has a home run}
\date{6 October 2026}

\begin{document}
\maketitle

\section*{Before you start}

This guide takes a computer with nothing on it and leaves you at the tool's own
instructions (the README's \emph{Verify} step). You do not need to be a programmer.
You will type commands into one black window; every command below can be copied and
pasted whole.

\textbf{What you are setting up.} The tool has three stages, and this guide installs
what each one needs, in that order:

\begin{center}
\begin{tabular}{@{}lp{0.44\textwidth}p{0.32\textwidth}@{}}
\toprule
\textbf{Stage} & \textbf{What it does} & \textbf{Needs} \\
\midrule
1. Scraping & collects an economy's laws from its official portals & Python packages + a browser engine \\
2. Extraction & splits the laws into numbered provisions & Python packages + Tesseract (for scans) \\
3. Mapping & maps each provision to the RDTII indicators & Python packages + Ollama (free engine) \\
\bottomrule
\end{tabular}
\end{center}

\textbf{Time and space.} About 30 minutes of typing and waiting, 2\,GB of disk.
A further 20 minutes and 9\,GB if you want the free local AI engine (Section 7).
Everything runs on your own computer; nothing is uploaded anywhere.

\textbf{The one window you need.}
\begin{itemize}[leftmargin=1.2em]
  \item \win: press the Start button, type \texttt{PowerShell}, open \emph{Windows PowerShell}.
  \item \mac: press \texttt{Cmd}+\texttt{Space}, type \texttt{Terminal}, press Enter.
\end{itemize}

\textbf{How to read this guide.} Each step first \emph{checks} whether something is
already on the computer. If the window answers with a version number, skip the install
and go to the next step. If it says \emph{not found}, \emph{not recognized} or
\emph{command not found}, run the install box underneath.

\section{Get the project onto the computer}

Download the project as a ZIP from the repository page on GitHub
(green \textbf{Code} button $\rightarrow$ \textbf{Download ZIP}), then unzip it to your
Desktop. GitHub names the unzipped folder after the branch or tag it came from:
\texttt{rdtii-rocky-next-main} or \texttt{rdtii-rocky-next-final-submission}. Rename it
to \texttt{rdtii-rocky-next} (right-click the folder, \emph{Rename}) so that the commands
below find it. You should end up with a folder called \texttt{rdtii-rocky-next} holding
folders named \texttt{interface}, \texttt{stages} and \texttt{submission}.

Now point the window at that folder. Every later command assumes you are inside it.

\win:
\begin{lstlisting}
cd $HOME\Desktop\rdtii-rocky-next
\end{lstlisting}

\mac:
\begin{lstlisting}
cd ~/Desktop/rdtii-rocky-next
\end{lstlisting}

Check you are in the right place --- this lists the project's files:

\begin{lstlisting}
dir        (Windows)
ls         (Mac)
\end{lstlisting}

\emph{If you prefer the developer's way:} if \texttt{git} is already installed, you can
instead run \texttt{git clone https://github.com/JohnChen-kmg/rdtii-rocky-next} followed
by \texttt{cd rdtii-rocky-next} and \texttt{git checkout final-submission}. The ZIP above
needs no extra program.

\section{Python: check, then install if missing}

Python is the language the tool is written in. Check first.

\win:
\begin{lstlisting}
python --version
\end{lstlisting}

\mac:
\begin{lstlisting}
python3 --version
\end{lstlisting}

\textbf{Good:} it answers \texttt{Python 3.12.x} or higher. Go to Section 3.

\textbf{Not good:} it says \texttt{not recognized} / \texttt{command not found}, it opens
the Microsoft Store, or the number is below 3.12. Install Python:

\begin{enumerate}[leftmargin=1.4em]
  \item Open \url{https://www.python.org/downloads/} and download the latest 3.12 or newer.
  \item \win: in the installer, \textbf{tick ``Add python.exe to PATH''} on the first screen
        before pressing Install. This one tick is what makes the command above work.
  \item \mac: open the downloaded \texttt{.pkg} and accept every default.
  \item \textbf{Close the window and open a new one} (it only notices the new Python then),
        go back to your project folder as in Section 1, and run the check again.
\end{enumerate}

\section{Give the project its own package area}

Rather than installing packages into your whole computer, the tool keeps them in one
folder inside the project, called \texttt{.venv}. This is one command and takes a few
seconds.

\win:
\begin{lstlisting}
python -m venv .venv
\end{lstlisting}

\mac:
\begin{lstlisting}
python3 -m venv .venv
\end{lstlisting}

Nothing is printed when it works. From here on, commands say \texttt{.venv\textbackslash Scripts\textbackslash}
(Windows) or \texttt{.venv/bin/} (Mac) so that the project's own Python is used.

\section{Stage 1 --- Scraping: the packages that fetch the laws}

\subsection*{Check}

\win:
\begin{lstlisting}
.venv\Scripts\python -c "import playwright, requests, tenacity, dateutil, yaml, jsonschema, cryptography; print('Stage 1 ready')"
\end{lstlisting}

\mac:
\begin{lstlisting}
.venv/bin/python -c "import playwright, requests, tenacity, dateutil, yaml, jsonschema, cryptography; print('Stage 1 ready')"
\end{lstlisting}

\textbf{Good:} it prints \texttt{Stage 1 ready}. \textbf{Not good:} it prints
\texttt{ModuleNotFoundError: No module named '...'}. Install below.

\subsection*{Install}

One command installs the packages of all three stages (they come from one list), so you
only run this once:

\win:
\begin{lstlisting}
.venv\Scripts\pip install -r requirements-demo.txt
\end{lstlisting}

\mac:
\begin{lstlisting}
.venv/bin/pip install -r requirements-demo.txt
\end{lstlisting}

It prints a long list and ends with \texttt{Successfully installed ...}. Allow 5 minutes.

Stage 1 also needs a browser engine to read portals that build their pages with
JavaScript. It is about 150\,MB:

\win:
\begin{lstlisting}
.venv\Scripts\python -m playwright install chromium
\end{lstlisting}

\mac:
\begin{lstlisting}
.venv/bin/python -m playwright install chromium
\end{lstlisting}

Then run the \emph{Check} box again; it should now print \texttt{Stage 1 ready}.

\subsection*{What these are for}

\begin{center}
\begin{tabular}{@{}ll@{}}
\toprule
\textbf{Package} & \textbf{Why Stage 1 needs it} \\
\midrule
\texttt{playwright} (+ Chromium) & drives a real browser for portals that need one \\
\texttt{requests}, \texttt{httpx} & downloads a document over the web \\
\texttt{tenacity} & waits and retries politely when a portal is busy \\
\texttt{python-dateutil} & reads the dates portals print in their own styles \\
\texttt{PyYAML} & reads the list of portals for each economy \\
\texttt{jsonschema} & checks the record of what was collected before handing it on \\
\texttt{cryptography} & Laws of Malaysia serves its listings encrypted \\
\bottomrule
\end{tabular}
\end{center}

\section{Stage 2 --- Extraction: reading documents, including scans}

\subsection*{Check}

\win:
\begin{lstlisting}
.venv\Scripts\python -c "import rdtii_p2.cli, pypdfium2, pytesseract, pydantic; print('Stage 2 ready')"
\end{lstlisting}

\mac:
\begin{lstlisting}
.venv/bin/python -c "import rdtii_p2.cli, pypdfium2, pytesseract, pydantic; print('Stage 2 ready')"
\end{lstlisting}

\subsection*{Install}

The extraction stage is part of the project, so it is installed from the project folder
itself:

\win:
\begin{lstlisting}
.venv\Scripts\pip install -e stages/p2-extract
\end{lstlisting}

\mac:
\begin{lstlisting}
.venv/bin/pip install -e stages/p2-extract
\end{lstlisting}

\subsection*{Tesseract --- only for scanned documents}

Lao PDR's laws are 96\% scans, so Stage 2 reads them with OCR. Check:

\begin{lstlisting}
tesseract --version
\end{lstlisting}

\textbf{Good:} it prints \texttt{tesseract 5.x}. \textbf{Not good:} \emph{not found}. Install:

\win:
\begin{lstlisting}
winget install UB-Mannheim.TesseractOCR
\end{lstlisting}

\mac{} (needs Homebrew; if \texttt{brew} is not found, install it from \url{https://brew.sh} first):
\begin{lstlisting}
brew install tesseract
\end{lstlisting}

Close the window and open a new one, then check again. The language packs (English, Lao,
Chinese, Portuguese) are already inside the project --- nothing to download.

\subsection*{What these are for}

\begin{center}
\begin{tabular}{@{}lp{0.5\textwidth}@{}}
\toprule
\textbf{Package} & \textbf{Why Stage 2 needs it} \\
\midrule
\texttt{pypdfium2} & opens a PDF and takes out its text, or a picture of the page \\
\texttt{pytesseract} + Tesseract & reads a scanned page into text (OCR) \\
\texttt{Pillow} & cleans up the page image before OCR \\
\texttt{beautifulsoup4}, \texttt{lxml}, \texttt{selectolax} & read a law published as a web page \\
\texttt{pydantic} & checks each provision has the fields the next stage expects \\
\texttt{jiwer} & measures how accurate the OCR was \\
\bottomrule
\end{tabular}
\end{center}

\section{Stage 3 --- Mapping: matching provisions to indicators}

\subsection*{Check}

\win:
\begin{lstlisting}
.venv\Scripts\python -c "import numpy, bm25s, pandas, openpyxl, anthropic, ollama; print('Stage 3 ready')"
\end{lstlisting}

\mac:
\begin{lstlisting}
.venv/bin/python -c "import numpy, bm25s, pandas, openpyxl, anthropic, ollama; print('Stage 3 ready')"
\end{lstlisting}

These came with the one \texttt{pip install} of Section 4. If anything is missing, run
that command again.

\subsection*{What these are for}

\begin{center}
\begin{tabular}{@{}ll@{}}
\toprule
\textbf{Package} & \textbf{Why Stage 3 needs it} \\
\midrule
\texttt{bm25s}, \texttt{PyStemmer} & finds the provisions worth reading, by keyword, at no cost \\
\texttt{numpy} & the arithmetic behind that search \\
\texttt{anthropic} & talks to the hosted engine (Engine A), if you choose it \\
\texttt{ollama} & talks to the free local engine (Engine B) \\
\texttt{pandas} & builds the per-economy reports \\
\texttt{openpyxl} & writes the Excel evidence workbook \\
\bottomrule
\end{tabular}
\end{center}

\section{The free local engine (optional, but it costs nothing)}

The tool declares two AI engines. \textbf{Engine A} is hosted (Claude) and needs a key you
pay for. \textbf{Engine B} runs on your own computer and is free. For Engine B, check:

\begin{lstlisting}
ollama --version
\end{lstlisting}

If it is not found, install Ollama from \url{https://ollama.com} (a normal installer for
Windows and Mac), then download the model the tool declares --- about 9\,GB, so allow time:

\begin{lstlisting}
ollama pull qwen2.5:14b
\end{lstlisting}

Confirm it is the exact model we declared (the digest should start \texttt{sha256:7cdf5a0187d5}):

\begin{lstlisting}
ollama list --digests
\end{lstlisting}

\textbf{You can skip this section} and still open the tool, read every filed result,
review rows and export them. You need an engine only to start a \emph{new} mapping run.

\section{Keys, settings and configuration files}

There is nothing to configure. No file to edit, no \texttt{.env}, no environment variable
to set: every location has a working default inside the project folder.

If you want to use the hosted engine (Engine A), you type its key \emph{into the tool} ---
Mapping tab, \textbf{Engine API keys} at the top of the page. The key is held in memory for
as long as the tool is open and is never written to disk, never shown again, and never sent
anywhere except to that provider.

\section{Start the tool}

\win: double-click \texttt{Start-RDTII-Rocky.cmd} in the project folder.\\
\mac: double-click \texttt{Start-RDTII-Rocky.command}.

The pages open in a window of their own, using the Edge or Chrome already on your machine.
Closing that window stops the tool.

A \mac{} may need one extra step, once: if the double-click opens the file as text, or
says it cannot be executed, the unzip has dropped its permission to run. Give it back from
the window, inside the project folder, then double-click again:

\begin{lstlisting}
chmod +x Start-RDTII-Rocky.command
\end{lstlisting}

If the double-click does nothing, start it from the window instead:

\win:
\begin{lstlisting}
.venv\Scripts\python interface/app.py
\end{lstlisting}

\mac:
\begin{lstlisting}
.venv/bin/python interface/app.py
\end{lstlisting}

\dots{} then open \url{http://127.0.0.1:8765/} in your browser.

\section{What a successful setup looks like}

You are ready for the tool's own instructions when all of this is true:

\begin{enumerate}[leftmargin=1.4em]
  \item The three check commands print \texttt{Stage 1 ready}, \texttt{Stage 2 ready},
        \texttt{Stage 3 ready}.
  \item The tool opens and shows the \textbf{Overview} page.
  \item Along the top of the page, the health dots read: \textbf{Stages} green;
        \textbf{Tesseract} green; \textbf{Chromium} green; \textbf{Ollama} green if you did
        Section 7, grey if you skipped it; \textbf{Key} grey until you type one, which is
        correct.
  \item Mapping tab $\rightarrow$ \textbf{3.3 Output} lists the finale's run (52 rows, six
        economies). Clicking a row shows the original law text beside its English, with the
        source link.
  \item Appendix tab $\rightarrow$ \textbf{Check the run layer} ends with
        \textbf{Finished}.
\end{enumerate}

If all five hold, go to the README's \emph{Verify} section and carry on there.

\section{If the window says something is not found}

\begin{center}
\begin{tabular}{@{}p{0.42\textwidth}p{0.52\textwidth}@{}}
\toprule
\textbf{What you see} & \textbf{What to do} \\
\midrule
\texttt{python is not recognized} (Windows) &
Python is missing or was installed without the PATH tick. Re-run the installer,
tick \emph{Add python.exe to PATH}, open a new window. \\[2pt]
Typing \texttt{python} opens the Microsoft Store &
Same answer: install from python.org with the PATH tick. \\[2pt]
\texttt{ModuleNotFoundError: No module named 'x'} &
A package is missing. Re-run \texttt{pip install -r requirements-demo.txt} from
Section 4, inside the project folder. \\[2pt]
\texttt{No such file or directory: requirements-demo.txt} &
You are in the wrong folder. Repeat the \texttt{cd} command of Section 1. \\[2pt]
\texttt{tesseract: command not found} &
Section 5. On Mac you need Homebrew first. Open a new window afterwards. \\[2pt]
The Ollama dot stays red &
Ollama is installed but not running: open the Ollama app once, then refresh the page. \\[2pt]
The page does not open &
Start the tool from the window (Section 9) --- the message there says what is wrong. \\[2pt]
A stage's \textbf{Start} button stays greyed out &
Press \textbf{Check} on that page first. The failed line tells you exactly what is missing. \\
\bottomrule
\end{tabular}
\end{center}

\vspace{4mm}
\noindent\rule{\textwidth}{0.4pt}

\small
Questions: John (Jiaxiang) Chen, \href{mailto:jiaxiangchen.kmg@gmail.com}{jiaxiangchen.kmg@gmail.com}.
The tool is Apache~2.0 licensed; this guide covers the release tagged \texttt{final-submission}.

\end{document}
